\documentclass[10pt]{article}
% Shared typographic and algebraic source.
\usepackage[letterpaper,margin=0.86in]{geometry}
\usepackage[T1]{fontenc}
\usepackage[utf8]{inputenc}
\usepackage{lmodern}
\usepackage{microtype}
\usepackage{amsmath,amssymb,amsthm,mathtools,mathrsfs}
\usepackage{booktabs,array,tabularx,longtable}
\usepackage{enumitem}
\usepackage[dvipsnames]{xcolor}
\usepackage{fancyhdr,titlesec}
\usepackage[hidelinks]{hyperref}

\definecolor{ink}{HTML}{1F2A33}
\definecolor{teal}{HTML}{226A73}
\definecolor{muted}{HTML}{65747E}
\definecolor{line}{HTML}{C5D1D3}
\definecolor{wash}{HTML}{F3F7F6}

\setlength{\parindent}{0pt}
\setlength{\parskip}{0.46em}
\setlength{\emergencystretch}{2em}
\allowdisplaybreaks[2]
\renewcommand{\arraystretch}{1.16}
\setlist{nosep,leftmargin=1.5em}

\pagestyle{fancy}
\fancyhf{}
\fancyhead[L]{\footnotesize\sffamily\color{muted}COPPALINI ARCHITECTURE}
\fancyhead[R]{\footnotesize\sffamily\color{muted}PHYSICAL REALITY INSIDE CONSTRAINT}
\fancyfoot[C]{\footnotesize\sffamily\color{muted}\thepage}
\renewcommand{\headrulewidth}{0.35pt}
\renewcommand{\headrule}{\hbox to\headwidth{\color{line}\leaders\hrule height \headrulewidth\hfill}}
\setlength{\headheight}{19pt}

\titleformat{\section}{\Large\bfseries\color{ink}}{\thesection.}{0.55em}{}
\titleformat{\subsection}{\large\bfseries\color{teal}}{\thesubsection.}{0.5em}{}
\titlespacing*{\section}{0pt}{1.42em}{0.48em}
\titlespacing*{\subsection}{0pt}{1.02em}{0.32em}

\newtheorem{theorem}{Theorem}
\newtheorem{proposition}[theorem]{Proposition}
\newtheorem{lemma}[theorem]{Lemma}
\theoremstyle{definition}
\newtheorem{definition}[theorem]{Definition}
\theoremstyle{remark}
\newtheorem{remark}[theorem]{Boundary}

\newcommand{\RR}{\mathbb R}
\newcommand{\CC}{\mathbb C}
\newcommand{\QQ}{\mathbb Q}
\newcommand{\ZZ}{\mathbb Z}
\newcommand{\Id}{I}
\newcommand{\tr}{\operatorname{tr}}
\newcommand{\im}{\operatorname{im}}
\newcommand{\kerop}{\operatorname{ker}}
\newcommand{\Span}{\operatorname{span}}
\newcommand{\diag}{\operatorname{diag}}
\newcommand{\snoc}{\operatorname{snoc}}
\newcommand{\Admit}{\operatorname{Admit}}
\newcommand{\Gate}{\operatorname{Gate}}
\newcommand{\Typed}{\operatorname{Typed}}
\newcommand{\Future}{\operatorname{future}}
\newcommand{\Scope}{\operatorname{sc}}
\newcommand{\SilScope}{\mathsf{Silence}_{\rm scope}}
\newcommand{\SilFit}{\mathsf{Silence}_{\rm fit}}

\newcommand{\papertitle}[2]{%
  \begin{center}
  {\color{teal}\rule{\textwidth}{1.35pt}}\\[0.38in]
  {\footnotesize\sffamily\bfseries\color{teal} COPPALINI ARCHITECTURE}\\[0.16in]
  {\fontsize{31}{36}\selectfont\bfseries\color{ink}#1\par}
  \vspace{0.15in}
  {\Large\color{teal}#2\par}
  \vspace{0.22in}
  {\color{line}\rule{\textwidth}{0.55pt}}\\[0.12in]
  {\normalsize Anthony Vito Coppa}\par
  {\small 9 November 2026}\par
  \vspace{0.30in}
  \end{center}
}

\newcommand{\sharedroot}{%
Let
\[
Q=\begin{pmatrix}1&1\\1&0\end{pmatrix},\qquad
\Omega=\begin{pmatrix}0&1\\-1&0\end{pmatrix},\qquad
C=-\Omega,
\]
\[
G=[Q,\Omega]=\begin{pmatrix}-2&1\\1&2\end{pmatrix},\qquad
H=\frac{G}{\sqrt5},\qquad
S=\frac{2Q-I}{\sqrt5}.
\]
The common exact identities are
\[
Q^2=Q+I,\qquad C^2=-I,\qquad G^2=5I,\qquad H^2=S^2=I,
\]
\[
H=CS,\qquad SG=\sqrt5\Omega=-GS,\qquad
L_s:=e^{sS},\qquad
L_s^{\mathsf T}GL_s=G,\qquad HL_sH=L_{-s}.
\]
}

\newcommand{\sharedoccurrence}[2]{%
A physical field has an interface
\[
\begin{aligned}
\mathfrak I_J:=(&J,\Gamma_J,\operatorname{Src}_J;
A_J,X_J,Y_J,T_J,\Sigma_J,\operatorname{Con}_J,\\
&\rho_J,O_J,K_J,B_J,\theta_J),
\qquad
\rho_J:X_J\to Y_J,\quad
O_J:X_J\to\mathcal O_J,\quad
\theta_J:X_J\to\{0,1\}.
\end{aligned}
\tag{#1}
\]
Its occurrence type is
\[
\begin{aligned}
\mathsf{Occ}_J:=\{&
(\mathfrak I_J;\Theta,x,\rho_J(x),\kappa,
\boldsymbol\varepsilon,\theta_J(x)):\\
&\Theta\in A_J,\ x\in X_J,\ \kappa\in K_J,\
\boldsymbol\varepsilon\in\RR_{\ge0}^{3}\}.
\end{aligned}
\tag{#2}
\]
An indexed occurrence $\mathsf{Occ}_{J,n}$ is an element of
$\mathsf{Occ}_J$.
}

\newcommand{\sharedmatching}{%
Let
\[
\mathcal A_8:=\{\infty\}\sqcup\ZZ_7,
\qquad
(\infty,0,1,2,3,4,5,6)
=(\mathsf{Gr},\mathsf{Li},\mathsf{Me},\mathsf{Ti},
\mathsf{Hz},\mathsf{Si},\mathsf{Sc},\mathsf{Mo}).
\]
For $r\in\ZZ_7$, set
\[
\mathsf M_r=\{\{\infty,r\}\}\cup
\{\{r+j,r-j\}:j=1,2,3\}.
\]
Then $E(K_8)=\bigsqcup_{r\in\ZZ_7}\mathsf M_r$, and the canonical matching is
\[
\mathsf M_0=\{\mathsf{Gravity\! -\! Light},\;
\mathsf{Memory\! -\! Motion},\;
\mathsf{Time\! -\! Scope},\;
\mathsf{Horizon\! -\! Silence}\}.
\]
}

\hypersetup{pdftitle={Memory Motion: Return, Mirror Transfer, and Relative Phase},pdfauthor={Anthony Vito Coppa},pdfsubject={Time-addressed transport, inverse return, and ordered physical occurrence}}
\begin{document}
\papertitle{Memory Motion}{Return, mirror transfer, and relative phase}

\begin{abstract}
Motion carries a physical state between time addresses. Memory is the inverse
condition by which a relation returns to its own address while the carriage
remains ordered. The axial metric flow advances a real coordinate by
$2\log\varphi$ and closes its quotient reading after that advance. The mirror
reverses the flow. Curvature records loop transport; material memory records
state update and readout. The equations distinguish a changing carrier state,
a returning read, and an ordered occurrence without erasing their difference.
\end{abstract}

\section{Motion and memory}

At time addresses $a\preceq b$, motion is a transport
\[
T_{b\leftarrow a}:X_a\longrightarrow X_b,
\qquad x_b=T_{b\leftarrow a}x_a.
\tag{MM.0}
\]
Time order composes exactly:
\[
T_{c\leftarrow b}T_{b\leftarrow a}=T_{c\leftarrow a},
\qquad T_{a\leftarrow a}=I_{X_a}.
\tag{MM.0a}
\]
Memory is the inverse condition of an invertible carriage:
\[
T_{a\leftarrow b}=T_{b\leftarrow a}^{-1},
\qquad T_{a\leftarrow b}T_{b\leftarrow a}=I_{X_a}.
\tag{MM.0b}
\]
Return is identity at its addressed carrier. A reader may close under a
transport while the state continues to differ at the local address.

\section{Axial lift and quotient}

\sharedroot

Let
\[
\varphi=\frac{1+\sqrt5}{2},\qquad
s_0=\log\varphi,\qquad
\varpi=2\log\varphi.
\tag{MM.1}
\]
The selected transport is
\[
L_{\rm ax}:=L_{s_0}=\frac{Q+2I}{\sqrt5},
\qquad T_{\rm ax}:=-L_{\rm ax}.
\tag{MM.2}
\]
Let $v_\pm$ denote the null eigenvectors of $Q$, with
$Qv_+=\varphi v_+$ and $Qv_-=-\varphi^{-1}v_-$. On
\[
X=\{z=av_++bv_-:ab\ne0\},
\tag{MM.3}
\]
define the lift coordinate
\[
\xi_{\rm lift}(z):=\log\left|\frac ab\right|.
\tag{MM.4}
\]

\begin{proposition}[Lift advance]
For $z\in X$,
\[
\xi_{\rm lift}(L_s z)=\xi_{\rm lift}(z)+2s,
\qquad
\xi_{\rm lift}(L_{\rm ax}z)=\xi_{\rm lift}(z)+\varpi.
\tag{MM.5}
\]
\end{proposition}

\begin{proof}
The flow $L_s=e^{sS}$ multiplies the two $S$-eigendirections by $e^s$ and
$e^{-s}$. Their coefficient ratio is therefore multiplied by $e^{2s}$.
Taking the logarithm proves the result.
\end{proof}

The axial reader is
\[
q_{\rm ax}:\RR\longrightarrow\RR/\varpi\ZZ,
\qquad q_{\rm ax}(\xi)=[\xi].
\tag{MM.6}
\]
It has the exact return-with-advance relation
\[
q_{\rm ax}\bigl(\xi_{\rm lift}(L_{\rm ax}z)\bigr)
=q_{\rm ax}(\xi_{\rm lift}(z)),
\qquad
\xi_{\rm lift}(L_{\rm ax}z)-\xi_{\rm lift}(z)=\varpi.
\tag{MM.7}
\]

\begin{theorem}[Quotient and lift]
The selected axial transport preserves the metric and reverses under the
mirror:
\[
L_{\rm ax}^{\mathsf T}GL_{\rm ax}=G,\qquad
HL_{\rm ax}H=L_{\rm ax}^{-1}.
\tag{MM.8}
\]
Its quotient return in MM.7 carries one reading across two distinct lift
coordinates.
\end{theorem}

\begin{proof}
The first two identities are instances of the shared root. The lift
difference in MM.7 is nonzero because $\varpi=2\log\varphi>0$. Hence the
states have distinct lift coordinates.
\end{proof}

\section{Reader, state, and record}

\sharedoccurrence{MM.9a}{MM.9b}

For $J=\mathsf{MM}$, let
\[
\operatorname{Ref}_{\rm MM}:\mathsf{Occ}_{\mathsf{MM}}
\to\mathcal R_{\rm MM},
\qquad
\Delta_{\rm MM}:\mathsf{Occ}_{\mathsf{MM}}
\to\mathcal D_{\rm MM}
\tag{MM.9c}
\]
be the declared reference and difference fields. The memory--motion receipt is
the projection
\[
r_n:=\operatorname{pr}_{\rm MM}(\mathsf{Occ}_{\mathsf{MM},n})
=(\operatorname{Ref}_{\rm MM}(\mathsf{Occ}_{\mathsf{MM},n}),
\Delta_{\rm MM}(\mathsf{Occ}_{\mathsf{MM},n}),y_n,\Theta_n).
\tag{MM.9d}
\]
An ordered record is
\[
R_{n+1}=\snoc(R_n,r_{n+1}).
\tag{MM.10}
\]
The record carries occurrence order while preserving every earlier occurrence
at its own address.

\begin{definition}[Reader-relative return]
For a transport $T$ and reader $\rho$, a return at $z$ is
\[
\rho(Tz)=\rho(z).
\tag{MM.11}
\]
A carrier difference is separately tested by $Tz\ne z$ or by a declared
difference map $\Delta(Tz,z)$.
\end{definition}

The axial instance is
\[
\rho=q_{\rm ax}\circ\xi_{\rm lift},\qquad T=L_{\rm ax},
\tag{MM.12}
\]
and satisfies both MM.11 and the nonzero lift condition in MM.7.

\section{Relative phase and mirror transfer}

Phase transport is
\[
U_\alpha=e^{\alpha C},
\]
with
\[
U_{\alpha+\beta}=U_\alpha U_\beta,\qquad
HU_\alpha H=U_{-\alpha}.
\tag{MM.13}
\]
For two addresses $a\preceq b$, their phase difference is
\[
\Delta\alpha_{b,a}:=\alpha(b)-\alpha(a)
\quad\text{in }\RR/2\pi\ZZ.
\tag{MM.14}
\]
The phase difference orders rotation between the two addresses. The axial lift
and the phase circle therefore remain separately readable coordinates.

A mirror-covariant reader obeys
\[
\rho(Hz)=\overline H_\rho(\rho(z)).
\tag{MM.15}
\]
When $\overline H_\rho$ is injective, equality and departure under $\rho$
transfer exactly between a trajectory and its mirror. For the projective
reader $\rho(z)=[z]$, the half-turn satisfies
\[
U_\pi z=-z\ne z,
\qquad [U_\pi z]=[z],
\tag{MM.15a}
\]
and gives an exact instance of reader-relative return.

\section{Curvature and physical memory}

For a connection $A$ on a bundle, curvature is the two-form
\[
F=dA+A\wedge A.
\tag{MM.16}
\]
A loop comparison requires a specified transport $P_\gamma$ and reader
$\rho_\gamma$. A reader-relative return takes the form
\[
\rho_\gamma(P_\gamma x)=\rho_\gamma(x).
\tag{MM.17}
\]
The loop relation carries a state around a closed physical path and compares
the returned read with its initial read.

A physical memory implementation likewise requires a material state space,
update, readout, and record:
\[
m_{n+1}=\mathcal U(m_n,u_n,d_n),
\qquad y_n=\rho_m(m_n).
\tag{MM.18}
\]
The state update, input, disturbance, readout, and occurrence address give a
material memory its physical form.

\end{document}
